\section{The revelations of the pulsing wave model}

It is clear that the macroscopic description of the system \ref{ch7:model:pulsing-waves}, \ie the equations \ref{ch7:EFFECTIVE-EQUATIONS}, are independent from certain mathematical objects of the microscopic model:

\begin{itemize}
	\item The effective description is not influenced by the constant deterministic velocity $\cdrift$ of the cAMP molecules.
	This might seem strange at first glance, but the actual drift for these particles is $\frac{\cdrift}{N}$, which goes to $0$ as $N$ grows to infinity: the result is expected.
	\item The system is independent from the choice of the chosen rotations $(\rotation_l)_{l=1}^\NumCampRelease$, meaning that even if every amoeba $\amba_t^p$ at every time $\btime_m$ shoots all its $\NumCampRelease$ children in the same random direction $\cgonia_{m,1}$ (\ie we set $\forall l. \rotation_l = \id$), the macroscopic description remains the same one.
	This is expected because what counts is that $\cgonia_{m,1}$ is uniformly distributed, and thus each direction $\rotation_l(\cgonia_{m,1})$ is in average equal to $0$, but it still remains a slightly counterintuitive fact about the system.
\end{itemize}

Then, the only parameters which do have a say also macroscopically are the number of children $\NumCampRelease$ at each birth time and the rate $\lambda$ with which such birth time are occuring.
With one simple idea for both parameters, we can say that the speed $\NumCampRelease \cdot \lambda$ at which the chemoattractant is released is telling us how much $\partial_t \adens$ is proportional to the member

\begin{equation*}
	\left[ \left( \nabla \potential
	\star \int_0^t \cdens_t^\tau 	d \tau
	\right)
	\adens
	\right]
\end{equation*}

The other thing which we can immediately observe, is the latter member we mentioned represents, in a certain way, a dependence of $\partial_t$ (the "present") from the past, \ie from all the states in $[0,t]$.
This happens because by integrating over $[0,t]$ in $\tau$ the expression $\cdens_t^\tau = \cdens^\tau (t,x) $ we are asking what happens at time $t$ to all the cAMP particles created along the time $[0,t]$.
This is even more evident if we substitute $\cdens_t^\tau$ with its analytical solution given by \ref{ch3:eq:eff-eq-camp-analytical-sol} and obtain \ref{ch7:PIDE}:

\begin{equation*}
	\partial_t \adens =  - \NumCampRelease\lambda \nabla \left[ \left( \nabla \int_0^t
	\potential \star \heatkernel_{t-\tau}^{\frac{d\sigma_\camp^2}{2}} \star \adens(\tau, \cdot) 
	d \tau
	\right)
	\adens
	\right]
	+ \frac{d\sigma_\amba^2}{2} \bigtriangleup \adens
\end{equation*}

Here $\adens$ appears also under the integral, and thus the memory effect (the direct dependence of the present $t$ from the events happening in $[0,t]$) is explicit.

\subsection{Comparison with the original paper of García}

The original paper who inspired this master thesis (\cite{ana}) modelled the movement of amoebas as particles attracting each others with a Coulomb-like potential (the cAMP molecules were not present in the microscopic IBM model), and found the famous Keller-Segel equation as an effective equation,

[[add the concept of IBM model in the introduction]]
[[add here up quote to the original Keller-Segel]]

in particular for the specific case in which the chemoattractant diffuses much faster than the organisms releasing it:

[[add here up quote to this specific Keller-Segel, always following Ana]]

\begin{equation}\label{conc:ANA-macro-amoebas}
	\left\{
	\begin{aligned}
		\partial_t \adens &= - \chi \nabla (\adens \nabla \cdens) + \bigtriangleup \adens
		\\
		-\bigtriangleup \cdens &= \adens
	\end{aligned}
	\right.
\end{equation}

or, solving $\cdens_t$ analytically:

\begin{equation}
		\partial_t \adens = - \chi \nabla ( (\cutoffSinAttraction \star \adens) \adens) + \bigtriangleup \adens
\end{equation}

where $\cutoffSinAttraction:= \frac{x}{2\pi \abs{x}^2}$ is the previously mentioned Coulomb-like attraction and $\chi$ is a positive constant called \emph{chemosensitivity}.
We modelled a slightly different scenario: the function $\attraction$ taken into consideration is Gaussian-like, and the diffusion coefficients $\diffusion_\amba, \diffusion_\camp$ of the amoebas and of the cAMP molecules are both finite (\ie of comparable order of magnitude).
We expect that it is possible to prove that the effective equation \ref{ch7:EFFECTIVE-EQUATIONS} is valid also in case of a Coulomb-like potential, so let us ignore the difference regarding $\attraction$ for now.
Then, to compare the two models, we want to ask what happens to \ref{ch7:PIDE} if we modulate the diffusions coefficient to be analogous to the ones treated in \cite{ana}, \ie when $\diffusion_\camp >> \diffusion_\amba:=\sqrt{\frac{2}{d}}$.
Consider the modified heat-kernel from \ref{ch3:eq:eff-eq-camp-analytical-sol}:

\begin{equation*}
	\heatkernel_s^{\frac{d\sigma_\camp^2}{2}}(x) = \frac{1}{
		\big( 2\pi d\sigma_\camp^2 s \big)^{\frac{d}{2}}
	}
	\exp \left(
	-\frac{x^2}{2 d\sigma_\camp^2 s}
	\right)
	\1_{s\geq 0} 
	= 
	\frac{d\: \Normal(0, d\sigma_\camp^2 s)}{d\:\lebsge} (x)
	\1_{s\geq 0} 
\end{equation*}

then in the limit case we do obtain

\begin{equation*}
		\heatkernel_s^{\frac{d\sigma_\camp^2}{2}}(x)
	\xrightarrow{\diffusion_\camp \rightarrow +\infty}
		\frac{d\: \Normal(0, "+\infty")}{d\:\lebsge} (x)
		\1_{s\geq 0} 
		\equiv
		0
\end{equation*}

and therefore, setting $\sigma_\amba := \sqrt{\frac{2}{d}}$, the effective equation \ref{ch7:EFFECTIVE-EQUATIONS} becomes

\begin{equation*}\label{conc:sigmaC-goes-infty}
	\partial_t \adens =  - \NumCampRelease\lambda \nabla \left[ \left( \nabla \int_0^t
	\potential \star 0 \star \adens(\tau, \cdot) 
	d \tau
	\right)
	\adens
	\right]
	+ \frac{d \frac{2}{d}}{2} \bigtriangleup \adens
	=
	\bigtriangleup \adens
\end{equation*}

which is a heat equation, which with the initial condition $\adens_0$ yields the explicit solution

\begin{equation*}
	\left(\heatkernel_t^1 \star \adens_0 \right) (x)
	= 
	\int 
	\frac{1}{
		\big( 2\pi)^{\frac{d}{2}}
	}
	\exp \left(
	-\frac{y^2}{2 t}
	\right)
	\adens_0 (x-y)
	dy
\end{equation*}

What is even more surprising is what happens in the other limit case, namely when $\sigma_\camp$ tends to $0^+$:

\begin{equation*}
	\heatkernel_s^{\frac{d\sigma_\camp^2}{2}}(x)
	\xrightarrow{\diffusion_\camp \rightarrow 0^+}
	\frac{d\: \Normal(0, 0)}{d\:\lebsge} (x)
	\1_{s\geq 0} 
	=
	\delta(x) \1_{s\geq 0}
\end{equation*}

In such case the macroscopic description of the amoebas' density becomes

\begin{equation}\label{conc:sigmaC-goes-0}
	\begin{aligned}
		\partial_t \adens &=  - \NumCampRelease\lambda \nabla \left[ \left( \nabla \int_0^t
		\potential \star 
		\delta(x) \1_{t-\tau\geq 0}
		\star \adens(\tau, \cdot) 
		d \tau
		\right)
		\adens
		\right]
		+
		\bigtriangleup \adens
		\\
		&=
		- \NumCampRelease\lambda \nabla \left[ \left( 
		\attraction
		\star 
		\int_0^t
		\adens(\tau, \cdot) 
		d \tau
		\right)
		\adens
		\right]
		+
		\bigtriangleup \adens
	\end{aligned}
\end{equation}

By setting $\chi:=\NumCampRelease\lambda$ we then obtain

\begin{equation*}\label{conc:sigmaC-goes-0}
	\partial_t \adens =
	- \chi \nabla \left[ \left( 
	\attraction
	\star 
	\int_0^t
	\adens(\tau, \cdot) 
	d \tau
	\right)
	\adens
	\right]
	+
	\bigtriangleup \adens
\end{equation*}

Now notice that the expression obtained is similar to the equation \ref{conc:ANA-macro-amoebas}: they differ in the convolution term, which is $(\cutoffSinAttraction \star \adens)$ in \ref{conc:ANA-macro-amoebas} and is $(\attraction \star \int_0^t \adens(\tau, \cdot) d\tau)$ in \ref{conc:sigmaC-goes-0}.
Though we can see that even if we would have the same attraction $\attraction$, the partial integro-differential equation (PIDE) \ref{conc:sigmaC-goes-0} still would have a memory of the past (the integration on $[0,t]$) that the Keller-Segel equation does not possess.

[[correct from here on. Besides, consider if to put such questions on the beginning of the conclusions]]

\section{Open questions on the topic}

\subsection{An interesting question: the more realistic mechanism of the chemical response}

[[write all this section properly]]









explain that other expansions of the model you have seen do not seem particularly interesting (eg L being a random parameter). Instead it is interesting to make the model realistic, having molecules which do answer to each others. Molecules which are able to speak chemically.

Build the Poisson process with the waiting times.
Then change these waiting times to some min(waiting time i, hitting time)

Explain that here the diffuculties is that your processes are then dependent on the system, and that you expect that the processes M t createdin this way to explode very fast towards infinity, or if you want that the frequency in time of birth times to get faster and faster.
Explain that this should make everything more diffucult to calculate.

\subsection{Technical "nice" conditions to assume or obtain for the pulsing wave models}

	\paragraph{Using a Coulomb-like potential}
	Again, as in the paper of Garcia (\cite{ana}), it would be interesting to use a Coulomb-like potential (analogous to the generalized attraction $\genAttraction$ used in the law of large numbers \ref{ch5:LLN}), and prove that the equation \ref{ch7:EFFECTIVE-EQUATIONS} still holds with $\attraction$ substituted with $\genAttraction$.
	As seen in other works previously (\cite{ana}, \cite{peter-boers}), the problem that arises is that $\genAttraction$ is not globally Lipschitz anymore, spoiling the same property for the drift $\drift(t, x)$.
	The solution is to show that $\drift(t, \cdot)$ is locally Lipschitz with a high probability (using essentially the same argument as the law of large numbers \ref{ch5:LLN}), and use such condition to bound the term $\drift(t, \sol_t) - \drift(t,\mfsol_t)$ which spawns in the proof of the final result \ref{ch7:final-result}.
	We expect this to be a not too difficult proof to build.

	\paragraph{Higher degree of convergence}
	As in the paper of Garcia (\cite{ana}) it would be interesting to obtain a higher degree of convergence for the distance between the microscopic and the mesoscopic description, namely going from the condition \ref{ch7:final-result} (for the basic model more precisely stated in \ref{ch6:final-result}) to a result of this kind:
	\begin{proposition}\label{conc:final-result-with-higher-convergence}
		Let $\sol_t^{(N)},\mfsol_t^{(N)}$ the solutions to the systems \ref{ch7:model:pulsing-waves} and \ref{ch7:MF-PARTICLES} respectively.
		Let $0\leq t \leq T$.
		Then, for each $\gamma>0$ and there is a positive, increasing and continuous function $\Fup_\gamma$ independent of $N$ such that
		\begin{equation}
			\PrbOf{
				\sup_{0 \leq t \leq T}
				\normp{\infty}{
					\sol_t^{(N)} - \mfsol_t^{(N)}
				} 
				\geq
				N^{-\cutparameter}
			}
			\leq\
			\Fup_{\gamma}(T)\
			N^{-\gamma}
		\end{equation}
	\end{proposition}
	Contrary to the usage of a Coulomb-like attraction, we expect this to be a tougher goal to pursue.	
	
\subsection{Numerics and simulations}
It could be an exciting task to search numerical solutions both for the PIDE and for the stochastic system \ref{ch7:model:pulsing-waves}, or to create simulations for the latter one.
Of course, we expect this not be easy because the PIDEs are generally considered very wild beasts, and because the system \ref{ch7:EFFECTIVE-EQUATIONS}, as underlined in chapter \ref{ch2}, is too stochastic to be precisely an SDE, therefore one could not directly apply the Euler-Maruyama method to it.
Still, it is not excluded that the author will one day set off on one of such adventures (keep an eye on this \href{https://github.com/carlolalu}{Github profile}).

\subsection{More technical and (probably) less interesting questions}

\paragraph{Higher dimensions}
We expect that if somebody would like to ask what happens with the Dictyostelium Discoideum amoebas moving in a $42$-dimensional space, the macroscopic description proved in this master thesis should still be valid, since along this work the dimension was a finite dimensional parameter $d$. 
We also expect this to be a futile question, unless somebody finds a $42$-dimensional natural phenomenon which can be modelled exactly in the same way.

\paragraph{The solution(s) to the PIDE and its properties}
As suggested in chapter \ref{ch3:prove-solutions?}, one might want to prove that the partial integro-differential equation \ref{ch7:PIDE} has a solution (with the required conditions) exploiting the Schauder fixed point theorem \ref{ch3:fixed-point-schauder-theorem}.
We expect this solution to exists, thus we regard this more as a technical verification rather than a quest bringing new insights.
One might also inquire about the properties of such $\adens(t, x)$, which would probably be a more interesting adventure, but since the object considered is a PIDE, we expect this not to be an easy task.

\paragraph{The macroscopic description of the cAMP molecules}
As suggested in \ref{ch3:macro-cAMP?}, one might inquire about proving formally that
\begin{equation*}
	"\1_{\hmactive_t \neq 0} \frac{1}{\hmactive_t}\sum_{m=1}^{\hmactive_t} \delta_{\camp_t^m} " =: \empmeas_t^{\camp,(N)} 
	\xrightarrow[\text{in law}]{N\longrightarrow +\infty} 
	\cdens(t,\cdot) :=
	\int_0^t \lambda \cdens_t^\tau 	d \tau
\end{equation*}
where the measure considered are treated as \rvs in the space of Borel probability measures with finite moments $\BorelPrbMeas_1$.